\subsubsection{Difference Quotients}\label{subsec:elliptic-regularity-difference-quotients}

Let $\Omega\subset\mathbb{R}^n$ and let $\{\bm{e}_1,\ldots,\bm{e}_n\}$ be the standard orthonormal basis for $\mathbb{R}^n$. For any function $u\in L^p(\Omega)$ we can formally define the difference quotient in the direction $\bm{e}_i$ to be
\begin{equation}\label{eq:coordinate-direction-difference-quotient-definition}
D_i^hu(\bm{x}):=\frac{u(\bm{x}+h\bm{e}_i)-u(\bm{x})}{h}.
\end{equation}
Of course, since $\bm{x}+h\bm{e}_i$ might extend beyond $\Omega$ for $\bm{x}$ near the boundary, this function might not be well defined for all $\bm{x}\in\Omega$. However, we can get the following result.

\begin{Lemma}
\label{lem:difference-quotient-weak-derivative-bound}
Let $u\in W^{1,p}(\Omega)$, $1\leq p\leq\infty$. Then for any $\Omega^{\prime}\subset\subset\Omega$ and any $h<\operatorname{dist}(\Omega^{\prime},\partial\Omega)$, we have $D_i^hu\in L^p(\Omega^{\prime})$ and
\begin{equation}\label{eq:difference-quotient-bound-by-weak-derivative}
\|D_i^hu\|_{L^p(\Omega^{\prime})}\leq\|u_{x_i}\|_{L^p(\Omega)}.
\end{equation}

\end{Lemma}

\begin{Proof}
Let $\Omega^{\prime}$ and $h$ satisfy the hypotheses of the lemma. For any $\xi\in[0,h]$ we define
\begin{equation}\label{eq:difference-quotient-translated-interior-domain}
\Omega_{\xi,i}^{\prime}:=\{\bm{x}\in\Omega\mid\bm{x}=\overline{\bm{x}}+\xi\bm{e}_i;\quad\overline{\bm{x}}\in\Omega^{\prime}\}.
\end{equation}
For $p\in[1,\infty)$ we first consider the case where $u\in C_b^1(\Omega)\cap W^{1,p}(\Omega)$. Then for any $\bm{x}\in\Omega^{\prime}$, we can use the fundamental theorem of calculus to write
\begin{equation}\label{eq:difference-quotient-average-derivative-identity}
\begin{aligned}
D_i^hu(\bm{x})&=\frac{u(\bm{x}+h\bm{e}_i)-u(\bm{x})}{h}\\
&=\frac{1}{h}\int_0^hu_{x_i}(\bm{x}+\xi\bm{e}_i)\,\mathrm{d}\xi.
\end{aligned}
\end{equation}
Thus, using Hölder\textquotesingle s inequality and switching orders of integration, we get\[\begin{aligned}
\int_{\Omega^{\prime}}|D_i^hu(\bm{x})|^p\,\mathrm{d}\bm{x}
&\leq\int_{\Omega^{\prime}}\frac{1}{h}\int_0^h|u_{x_i}(\bm{x}+\xi\bm{e}_i)|^p\,\mathrm{d}\xi\,\mathrm{d}\bm{x}\\
&\leq\frac{1}{h}\int_0^h\int_{\Omega_{\xi,i}^{\prime}}|u_{x_i}(\bm{x})|^p\,\mathrm{d}\bm{x}\,\mathrm{d}\xi\\
&\leq\int_\Omega|u_{x_i}|^p\,\mathrm{d}\bm{x}.
\end{aligned}\]By Theorem~\ref{thm:sobolev-interior-smooth-density}, this inequality extends to the whole space by taking limits.

For $p=\infty$ we note that \eqref{eq:difference-quotient-average-derivative-identity} holds in the sense of distributions. Since the test functions are dense in $L^1$, the inequality follows from manipulating and estimating $D_i^hu$ and $u_{x_i}$ as linear functionals on sets of test functions. We leave this to the reader.
\end{Proof}
